% Slide 12 — Ambiente de simulação e métricas
\begin{frame}{Módulo 3 -- Ambiente de simulação}
  \begin{columns}[T,onlytextwidth]
    \column{0.5\textwidth}
      \textbf{Parâmetros fixos}
      \small
      \begin{itemize}
        \item Horizonte $T=38$ ciclos comerciais
        \item Prazo de entrega $L=2$ ciclos
        \item $h=1{,}0$; $c_s=5{,}0$; $c_o^{\text{fix}}=50{,}0$; $c_o^{\text{var}}=0{,}5$
        \item Restrição de viabilidade: NS $\geq \alpha_{\min}=0{,}70$
      \end{itemize}
    \column{0.5\textwidth}
      \textbf{Indicadores (KPIs)}
      \small
      \begin{itemize}
        \item \textbf{CTI}: custo total de inventário
        \item \textbf{NS}: nível de serviço
        \item \textbf{TR}: taxa de ruptura
        \item \textbf{BE}: efeito \textit{bullwhip} ($\mathrm{Var}(Q)/\mathrm{Var}(D)$)
        \item \textbf{FP}: frequência de pedidos
      \end{itemize}
  \end{columns}
  \vfill
  \small Todas as políticas são avaliadas no \highlight{mesmo ambiente}: mesma
  demanda real, mesmos custos, mesmas sementes e os mesmos KPIs (CTI, NS, TR,
  BE e FP).
  \note{
    O ambiente de simulação é o que torna a comparação entre políticas
    válida: todas rodam sob as mesmas séries de demanda real, os mesmos
    custos e o mesmo prazo de entrega de dois ciclos. Cinco indicadores são
    calculados ao final de cada simulação. O CTI é o critério principal de
    custo; o NS funciona como restrição de viabilidade -- políticas com
    nível de serviço abaixo de 0,70 não são consideradas dominantes mesmo
    que tenham custo baixo; TR, BE e FP são complementares, com o BE sendo
    usado mais adiante na análise de estabilidade.
  }
\end{frame}
